\section{Introduction}
\label{sec:introduction}

For integers \(s,t\ge2\), the Ramsey number \(r(s,t)\) is the least
integer \(n\) such that every simple graph on \(n\) vertices contains
either a complete graph \(K_s\) or an independent set of size \(t\).
We determine the logarithmic exponent in the off-diagonal case \(s=5\).
We write \(\alpha(G)\) for the largest size of an independent set in \(G\).
All logarithms are natural.

\begin{theorem}
\label{thm:main}
There is an absolute constant \(C>0\) such that, for every
\(\varepsilon>0\), all sufficiently large integers \(t\) satisfy
\[
 \frac{t^4}{(\log t)^{3+\varepsilon}}
 \ \le\ r(5,t)\ \le\
 C\frac{t^4}{(\log t)^3}.
\]
The threshold for \(t\) may depend on \(\varepsilon\).
Consequently,
\[
 \lim_{t\to\infty}
 \frac{4\log t-\log r(5,t)}{\log\log t}=3.
\]
\end{theorem}

The theorem resolves the sharp logarithmic-exponent question for
\(r(5,t)\). It leaves open a constant-factor asymptotic for
